% Section 6: CLASS REFERENCE

\section{Class Reference}

This section provides the API reference for each class in the VentCon2 system, including class structures, key attributes, and method signatures.

\subsection{Class Summary Tables}
The following tables document each class with its attributes and methods.

\begin{table}[H]
\centering
\caption{SettingsHandler Class}
\label{tab:settings-handler}
\begin{tabular}{|l|l|p{6cm}|}
\hline
\textbf{Type} & \textbf{Member} & \textbf{Description} \\
\hline
\multicolumn{3}{|c|}{\textit{Attributes}} \\
\hline
double & Kp, Ki, Kd & PID controller gains \\
double & setpoint & Target pressure (bar) \\
int & pwm\_freq, pwm\_res & PWM configuration \\
int & pid\_sample\_time, control\_freq\_hz & PID sample time and control loop rate \\
float & filter\_strength & Low-pass filter coefficient \\
bool & antiWindup, hysteresis & Control enhancements \\
float & hystAmount & Hysteresis compensation (\%) \\
SliderLimits & sp\_, kp\_, ki\_, kd\_limits & UI slider limits \\
float & sensor\_min/max\_pressure, voltage & Sensor calibration \\
ChartSettings & chart\_settings & Live chart axes \\
\hline
\multicolumn{3}{|c|}{\textit{Methods}} \\
\hline
bool & load() & Load settings from LittleFS \\
bool & save() & Save settings to LittleFS \\
void & resetToDefaults() & Restore factory defaults \\
void & update(fn) & Run fn under the mutex \\
void & withLock(fn) & Run fn under the mutex \\
void & snapshotTo(SettingsSnapshot\&) & Atomic multi-field read \\
void & printStoredSettings() & Debug output to Serial \\
\hline
\end{tabular}
\end{table}

\begin{table}[H]
\centering
\caption{SensorManager Class}
\label{tab:sensor-manager}
\begin{tabular}{|l|l|p{6cm}|}
\hline
\textbf{Type} & \textbf{Member} & \textbf{Description} \\
\hline
\multicolumn{3}{|c|}{\textit{Attributes}} \\
\hline
Adafruit\_ADS1015 & ads & External ADC instance \\
bool & ads\_found & ADC detection flag \\
int16\_t & adc\_value & Raw ADC reading \\
float & voltage, raw\_pressure & Converted values \\
float & filtered\_pressure & After low-pass filter \\
\hline
\multicolumn{3}{|c|}{\textit{Methods}} \\
\hline
bool & initialize() & Setup I2C and ADC \\
void & readSensor() & Read and convert ADC \\
float & getPressure() & Get filtered pressure \\
float & getRawPressure() & Get unfiltered pressure \\
float & getVoltage() & Get sensor voltage \\
int16\_t & getADCValue() & Raw ADC reading \\
bool & isADSFound() & ADS1015 detected? \\
float* & getLastFilteredPressurePtr() & Pointer for web interface \\
bool & isPressureSafe() & Over-pressure safety check \\
void & resetFilterState() & Reset filter state \\
\hline
\end{tabular}
\end{table}

\begin{table}[H]
\centering
\caption{WebHandler Class}
\label{tab:web-handler}
\begin{tabular}{|l|l|p{6cm}|}
\hline
\textbf{Type} & \textbf{Member} & \textbf{Description} \\
\hline
\multicolumn{3}{|c|}{\textit{Attributes}} \\
\hline
WebServer & webServer & ESP32 HTTP server (port 80) \\
bool & webServerEnabled & Server state flag \\
int & connectedClients & Active client counter \\
String[2] & connectedMACs & Connected client MAC addresses \\
\hline
\multicolumn{3}{|c|}{\textit{Methods}} \\
\hline
void & handleRoot() & Stream HTML from PROGMEM \\
void & handleValues() & JSON API for live data \\
void & handleSet() & Process parameter changes \\
void & handleResetPID() & Reset PID state \\
void & handleSliderLimits() & GET/POST slider limits \\
void & handleChartSettings() & GET/POST chart axes \\
void & initializeWiFiAP() & Start AP with WIFI\_AP mode \\
void & scanWiFiNetworks() & Return JSON network list \\
\hline
\end{tabular}
\end{table}

\begin{table}[H]
\centering
\caption{ControlSystem Class}
\label{tab:control-system}
\begin{tabular}{|l|l|p{6cm}|}
\hline
\textbf{Type} & \textbf{Member} & \textbf{Description} \\
\hline
\multicolumn{3}{|c|}{\textit{Attributes}} \\
\hline
PID* & pid & PID controller instance \\
double* & pressureInput & Pointer to sensor reading \\
double* & pwmPIDoutput & Pointer to valve output \\
uint32\_t* & actualPwm & Mapped valve PWM value \\
bool* & manualPWMMode & Manual override flag \\
\hline
\multicolumn{3}{|c|}{\textit{Methods}} \\
\hline
void & initializeControlSystem() & Setup PID with settings \\
void & processControlLoop() & Execute one control cycle \\
uint32\_t & mapPIDoutputToPwmValve() & Map PID output to valve duty \\
void & applyHysteresisCompensation() & Hysteresis offset \\
void & handleEmergencyShutdown() & Safe state on error \\
\hline
\end{tabular}
\end{table}

\begin{table}[H]
\centering
\caption{AutoTuner Class}
\label{tab:auto-tuner}
\begin{tabular}{|l|l|p{6cm}|}
\hline
\textbf{Type} & \textbf{Member} & \textbf{Description} \\
\hline
\multicolumn{3}{|c|}{\textit{Attributes}} \\
\hline
bool & autoTuneRunning & Tuning active flag \\
float & autoTuneSetpoint & Target for oscillation \\
int & currentCycle & Oscillation counter \\
ulong[] & cycleTimes & Period measurements \\
float[] & cycleAmplitudes & Amplitude measurements \\
enum & currentTuningRule & Algorithm selection \\
\hline
\multicolumn{3}{|c|}{\textit{Methods}} \\
\hline
void & start() & Begin relay tuning \\
void & stop() & Abort and restore PID \\
void & process() & Execute tuning step \\
float & getProgress() & 0.0--1.0 completion \\
void & setTuningRule() & Select calculation method \\
\hline
\end{tabular}
\end{table}

\begin{table}[H]
\centering
\caption{CommandProcessor Class}
\label{tab:command-processor}
\begin{tabular}{|l|l|p{6cm}|}
\hline
\textbf{Type} & \textbf{Member} & \textbf{Description} \\
\hline
\multicolumn{3}{|c|}{\textit{Attributes}} \\
\hline
String & commandBuffer & Incoming serial data \\
Pointers & settings, sensor, etc. & References to all subsystems \\
\hline
\multicolumn{3}{|c|}{\textit{Methods}} \\
\hline
void & handleSerialInput() & Read and buffer chars \\
bool & executeCommand() & Parse and dispatch \\
void & showHelp() & Print command list \\
void & showStatus() & Print system state \\
void & handlePIDCmd() & PID parameter commands \\
void & handlePWMCmd() & Manual PWM commands \\
\hline
\end{tabular}
\end{table}

\begin{table}[H]
\centering
\caption{TaskManager Class}
\label{tab:task-manager}
\begin{tabular}{|l|l|p{6cm}|}
\hline
\textbf{Type} & \textbf{Member} & \textbf{Description} \\
\hline
\multicolumn{3}{|c|}{\textit{Attributes}} \\
\hline
TaskHandle\_t & controlTaskHandle & FreeRTOS control task \\
TaskHandle\_t & networkTaskHandle & FreeRTOS network task \\
bool & taskRunning & Tasks active flag \\
\hline
\multicolumn{3}{|c|}{\textit{Methods}} \\
\hline
bool & createTasks() & Spawn FreeRTOS tasks \\
void & startTasks() & Resume suspended tasks \\
void & stopTasks() & Suspend running tasks \\
bool & getTaskStatus() & Query task states \\
void & monitorTasks() & Watchdog and recovery \\
\hline
\end{tabular}
\end{table}

\begin{table}[H]
\centering
\caption{Constants.h Namespaces}
\label{tab:constants}
\begin{tabular}{|l|l|p{6cm}|}
\hline
\textbf{Namespace} & \textbf{Constant} & \textbf{Value / Description} \\
\hline
NetworkConfig & AP\_SSID & "VENTCON\_AP" \\
NetworkConfig & AP\_PASSWORD & "ventcon12!" (WPA2-PSK) \\
HardwareConfig & SOLENOID\_PIN & GPIO for valve PWM \\
HardwareConfig & ANALOG\_PRESS\_PIN & ADC input pin \\
SensorConfig & SENSOR\_MIN\_BAR & 0.0 (minimum range) \\
SensorConfig & SENSOR\_MAX\_BAR & 10.0 (maximum range) \\
ValveConfig & VALVE\_MIN\_DUTY & 20.0\% (minimum effective duty cycle) \\
ValveConfig & VALVE\_MAX\_DUTY & 90.0\% (maximum safe duty cycle) \\
OperatingDefaults & HARD\_SP\_MAX\_BAR & 10.0 (firmware setpoint ceiling) \\
SensorConfigDefaults & SENSOR\_MIN/MAX\_VOLTAGE & 0.5 / 4.5 V \\
SettingsDefaults & PID\_SAMPLE\_TIME & 10 ms \\
SettingsDefaults & CONTROL\_FREQ\_HZ & 1000 Hz \\
\hline
\end{tabular}
\end{table}

\subsection{SettingsHandler Class}

\subsubsection{Purpose}
The SettingsHandler class encapsulates all system configuration parameters and provides persistent storage capabilities using LittleFS. It serves as the central configuration repository for all system components.

\subsubsection{Key Responsibilities}
\begin{itemize}
    \item Manage PID controller parameters (Kp, Ki, Kd, setpoint)
    \item Store PWM configuration (frequency, resolution)
    \item Handle sensor filtering parameters
    \item Provide JSON serialization/deserialization
    \item Automatic load/save to flash storage
    \item Parameter validation and constraints
\end{itemize}

\subsubsection{Key Attributes}
\begin{lstlisting}[caption=SettingsHandler Class Key Members - Actual Implementation]
// SliderLimits structure for UI customization
struct SliderLimits {
    float min;                    // Minimum slider value
    float max;                    // Maximum slider value
    float step;                   // Slider step increment
};

class SettingsHandler {
public:
    // PID Parameters
    double Kp;                    // Proportional gain
    double Ki;                    // Integral gain  
    double Kd;                    // Derivative gain
    
    // System Parameters
    float filter_strength;        // Low-pass filter coefficient
    double setpoint;              // Target pressure in bar
    int pwm_freq;                 // PWM frequency in Hz
    int pwm_res;                  // PWM resolution in bits
    int pid_sample_time;          // PID sample time in ms
    int control_freq_hz;          // Control loop frequency
    
    // Advanced Features
    bool antiWindup;              // Anti-windup enable flag
    bool hysteresis;              // Hysteresis compensation
    float hystAmount;             // Hysteresis amount (percentage)
    
    // Slider Limits (user-configurable via web UI)
    SliderLimits sp_limits;       // Setpoint slider limits
    SliderLimits kp_limits;       // Kp slider limits
    SliderLimits ki_limits;       // Ki slider limits
    SliderLimits kd_limits;       // Kd slider limits
    
    // Sensor Calibration (user-configurable pressure/voltage mapping)
    float sensor_min_pressure;    // Minimum pressure (bar)
    float sensor_max_pressure;    // Maximum pressure (bar)
    float sensor_min_voltage;     // Sensor voltage at min pressure (V)
    float sensor_max_voltage;     // Sensor voltage at max pressure (V)
    
    // Chart Axis Settings
    ChartSettings chart_settings; // Live chart axes and time window
    
    // Constructor with default values
    SettingsHandler();
    
    // Methods
    bool load();                  // Load from LittleFS
    bool save();                  // Save to LittleFS
    void resetToDefaults();       // Reset to default values
    void printStoredSettings();   // Display stored settings
    
    // Mutex-protected access
    template<typename Fn> void update(Fn fn);   // Run fn(*this) under mutex
    template<typename Fn> void withLock(Fn fn); // Run fn() under mutex
    void snapshotTo(SettingsSnapshot& out);     // Atomic multi-field read
};
\end{lstlisting}

\subsection{SensorManager Class}

\subsubsection{Purpose}
The SensorManager class handles all sensor-related operations, providing a unified interface for pressure sensing with automatic fallback mechanisms and signal processing.

\subsubsection{Key Responsibilities}
\begin{itemize}
    \item ADS1015 external ADC initialization and communication
    \item ESP32 internal ADC fallback when external ADC fails
    \item Voltage to pressure conversion calculations
    \item Low-pass filtering for noise reduction
    \item Sensor health monitoring and diagnostics
\end{itemize}

\subsubsection{Key Attributes}
\begin{lstlisting}[caption=SensorManager Class Structure]
class SensorManager {
private:
    Adafruit_ADS1015 ads;         // External ADC instance
    bool ads_found;               // ADC availability flag
    int16_t adc_value;           // Raw ADC reading
    float voltage;               // Converted voltage
    float raw_pressure;          // Unfiltered pressure
    float filtered_pressure;     // Filtered pressure
    SettingsHandler* settings;          // Configuration reference
    
public:
    SensorManager(SettingsHandler* settings);
    bool initialize();           // Hardware initialization
    void readSensor();          // Main sensor reading function
    float getPressure();        // Get filtered pressure
    bool isADSFound();         // Check ADC status
    float* getLastFilteredPressurePtr(); // For web interface
    bool isPressureSafe();     // Over-pressure safety check
    void resetFilterState();   // Reset filter state
};
\end{lstlisting}

\subsubsection{Sensor Specifications}
\begin{table}[H]
\centering
\begin{tabular}{@{}ll@{}}
\toprule
\textbf{Parameter} & \textbf{Value} \\
\midrule
Voltage Range & 0.5V - 4.5V \\
Pressure Range & 0 - 10 bar \\
ADC Resolution & 12-bit (ADS1015) \\
Sample Rate & Up to 1600 SPS \\
I2C Address & 0x48 \\
Fallback Pin & A0 (ESP32 internal) \\
\bottomrule
\end{tabular}
\caption{Sensor Configuration Parameters}
\end{table}

\subsection{AutoTuner Class}

\subsubsection{Purpose}
The AutoTuner class implements relay-based auto-tuning for PID controller parameters using the Ziegler-Nichols frequency response method and other tuning algorithms. The tuning algorithm and rules are described in Section~\ref{sec:auto-tuning}.

\subsubsection{Class Structure}
\begin{lstlisting}[caption=AutoTuner Class Key Members - Actual Implementation]
class AutoTuner {
private:
    // Auto-tuning state variables
    bool autoTuneRunning;
    unsigned long autoTuneStartTime;
    unsigned long lastTransitionTime;
    float autoTuneOutputValue;
    float autoTuneSetpoint;
    bool autoTuneState;
    int currentCycle;
    
    // Cycle data collection
    static constexpr int MAX_CYCLES = 20;
    unsigned long cycleTimes[MAX_CYCLES];
    float cycleAmplitudes[MAX_CYCLES];
    
    // Amplitude tracking
    float maxPressure;
    float minPressure;
    bool firstCycleComplete;
    
    // Configuration parameters
    float testSetpoint;
    float minPwmValue;
    float maxPwmValue;
    unsigned long minCycleTime;
    TuningRule currentTuningRule;
    float tuningAggressiveness;
    
    // System references
    SettingsHandler* settings;
    PID* pid;
    double* pressureInput;
    int* pwmMaxValue;
    
public:
    AutoTuner(SettingsHandler* settings, PID* pid, double* pressureInput, int* pwmMaxValue);
    void start();
    void stop(bool calculateParameters = false);
    void process();
    bool isRunning() const;
    float getOutputValue() const;
    void acceptParameters();
    void rejectParameters();
    void setTestSetpoint(float setpoint);
    void setTuningRule(TuningRule rule);
    void setAggressiveness(float aggr);
    void setMinMaxPWM(float minPwm, float maxPwm);
    void setMinCycleTime(unsigned long cycleTime);
    float getTestSetpoint() const;
    TuningRule getTuningRule() const;
    float getAggressiveness() const;
    float getMinPWM() const;
    float getMaxPWM() const;
    float getEffectiveAmplitude() const;
    unsigned long getMinCycleTime() const;
    void printTuningRules() const;
    void printConfiguration() const;
};
\end{lstlisting}

\subsection{ControlSystem Class}

\subsubsection{Purpose}
The ControlSystem class implements the main control loop, integrating PID control, sensor reading, valve control, and safety mechanisms in a real-time task. The control algorithms (valve mapping, anti-windup, hysteresis) are described in Section~\ref{sec:control-theory}.

\subsubsection{Control Loop Structure}
\begin{lstlisting}[caption=Control Loop Implementation - Actual Code from ControlSystem.cpp]
void ControlSystem::processControlLoop() {
    TickType_t lastWakeTime = xTaskGetTickCount();
    
    // Use configurable frequency
    TickType_t frequency = pdMS_TO_TICKS(1000 / settings->control_freq_hz);
    
    static unsigned long lastCycleEnd = 0;
    
    while(true) {
        // Update frequency if settings changed
        frequency = pdMS_TO_TICKS(1000 / settings->control_freq_hz);
        unsigned long taskStartTime = micros();
        
        // ====== Sensor Reading using SensorManager ======
        sensorManager->readSensor();
        
        // Save last pressure before updating, used for hysteresis compensation
        lastPressure = *pressureInput;
        
        // Update Input value for PID before computation
        *pressureInput = sensorManager->getPressure();
        
        // ====== Analog Pressure Signal Output =====
        handleAnalogPressureOutput();
        
        // Process auto-tuning if active
        if (autoTuner && autoTuner->isRunning()) {
            autoTuner->process();
        }
        // PID calculation and PWM output (only if not in auto-tune mode)
        else if (!*manualPWMMode) {
            // Store previous output for anti-windup check
            double previousOutput = *pwmPIDoutput;
            
            // Constrain output to valid range
            *pwmPIDoutput = constrain(*pwmPIDoutput, 0, *pwmFullScaleRaw); 

            // Compute PID output
            pid->Compute();
            
            // Anti-windup for deadband and saturation
            if (settings->antiWindup) {
                float pidPercent = (*pwmPIDoutput / *pwmFullScaleRaw) * 100.0;
                
                if ((pidPercent < ValveConfig::VALVE_MIN_DUTY && *pwmPIDoutput < previousOutput) ||
                    (pidPercent > ValveConfig::VALVE_MAX_DUTY && *pwmPIDoutput > previousOutput)) {
                    // Reset the PID to prevent integral accumulation
                    pid->SetMode(PID::Manual);
                    pid->SetMode(PID::Automatic);
                }
            }
            
            // Update PWM output
            updatePWMOutput();
        }
        
        // ====== Continuous Data Output (Serial) ======
        handleContinuousDataOutput();
        
        // ====== Task Timing Management ======
        unsigned long taskEndTime = micros();
        lastCycleTime = taskEndTime - taskStartTime;
        
        vTaskDelayUntil(&lastWakeTime, frequency);
    }
}
\end{lstlisting}

\subsection{WebHandler Class}

\subsubsection{Purpose}
The WebHandler class provides a comprehensive web interface for system monitoring and control, implementing a WiFi Access Point with mDNS and HTTP endpoints. It serves as the primary user interface for remote monitoring and real-time parameter adjustment of the pressure control system. The web interface architecture is described in Section~\ref{sec:web-interface}.

\subsubsection{Network Configuration}
\begin{itemize}
    \item \textbf{SSID:} VENTCON\_AP
    \item \textbf{Password:} ventcon12!
    \item \textbf{IP Address:} 192.168.4.1
    \item \textbf{DNS:} mDNS (http://ventcon.local)
    \item \textbf{Max Clients:} 2 simultaneous connections
    \item \textbf{Channel:} Configurable (use WIFI CHANNEL command)
    \item \textbf{Security:} WPA2-PSK encryption
\end{itemize}

\subsubsection{Class Structure}
\begin{lstlisting}[caption=WebHandler Class - Actual Implementation]
class WebHandler {
private:
    static WebHandler* instance;  // For WiFi event callbacks
    
    // Owned network components
    WebServer webServer;
    
    // References to system components (dependency injection)
    SettingsHandler* settings;
    PID* pid;
    double* pressureInput;
    double* pwmPIDoutput;
    bool* ads_found;
    int* pwmFullScaleRaw;
    float* last_filtered_pressure;
    
    // WiFi connection management
    int connectedClients;
    String connectedMACs[NetworkConfig::MAX_CLIENTS];
    bool webServerEnabled;
    
    // WiFi AP configuration
    const char* ap_ssid;
    const char* ap_password;
    IPAddress ap_ip;       // 192.168.4.1
    IPAddress ap_gateway;  // 192.168.4.1
    IPAddress ap_subnet;   // 255.255.255.0
    
    // Private helper methods
    String getContentType(String filename);
    bool handleFileRead(String path);  // Chunked file serving
    
    // Route handlers
    void handleRoot();      // Memory-efficient streaming
    void handleSet();
    void handleValues();
    void handleResetPID();
    void handleSliderLimits();  // GET/POST slider limits API
    void handleChartSettings(); // GET/POST chart axis settings API
    
public:
    WebHandler(SettingsHandler* settings, PID* pid, ...);
    ~WebHandler();
    
    void setupRoutes();
    void updatePWM();
    void scanWiFiNetworks();
    void changeWiFiChannel(int channel);
    void setupWiFiEvents();
    void initializeWiFiAP();
    
    // Getters
    int getConnectedClients() const;
    bool isWebServerEnabled() const;
    int getWiFiChannel() const;
};
\end{lstlisting}

\subsection{CommandProcessor Class}

\subsubsection{Purpose}
The CommandProcessor class provides a comprehensive serial command interface for advanced system control, debugging, and configuration. All commands are case-insensitive and accessible through the serial console at 115200 baud. The full command reference is provided in Section~\ref{sec:serial-commands}.

\subsection{TaskManager Class}

\subsubsection{Purpose}
The TaskManager class orchestrates FreeRTOS tasks for optimal real-time performance, separating control operations from network operations across ESP32 cores. Task distribution details are described in Section~\ref{sec:software-arch}.

\subsection{Logger Class}

\subsubsection{Purpose}
The Logger class provides centralized, rate-limited serial logging with per-category rate limiting and compile-time level filtering. It is initialized once in \texttt{main.cpp} via \texttt{Logger::init()} and used across all components.

\subsubsection{Key Members}
\begin{lstlisting}[caption=Logger Class - Key Members]
class Logger {
public:
    static void init();                       // Initialize the logger
    static void setEnabled(bool en);          // Enable/disable output
    static bool isEnabled();                  // Check if enabled
    static void setRateLimit(unsigned long ms); // Per-category rate limit
    static bool shouldLog(LogCategory cat);   // Rate-limit check
    static const char* getCategoryName(LogCategory cat);
};

enum LogLevel { LOG_NONE = 0, LOG_ERROR = 1, LOG_WARN = 2, LOG_INFO = 3, LOG_DEBUG = 4 };
enum LogCategory { CAT_CONTROL, CAT_NETWORK, CAT_SENSOR, CAT_AUTOTUNE, CAT_SYSTEM, CAT_COUNT };
\end{lstlisting}
